\section{Extended methods}
\label{sec:si:methods}

This section states in full the procedures the article summarises. Nothing here is a different
analysis from the one reported there; it is the detail a reader would need in order to reproduce
it without reading the code.

\subsection{The three clusterless generators}
\label{sec:si:generators}

Let $X \in \mathbb{R}^{n \times d}$ be the standardised feature matrix for a scope, with $n$
player-seasons and $d = \ArcFeatures{}$ features. Let $s_k(X)$ be the mean silhouette of the
partition $k$-means returns at $k$ clusters, computed on a fixed subsample of at most four
thousand rows when $n$ exceeds that, since the silhouette is quadratic in $n$.

The Gaussian generator draws $n$ rows from $\mathcal{N}(\hat\mu, \hat\Sigma)$ with the empirical
mean and covariance of $X$. It matches the data in size, dimension and covariance and in nothing
else: every marginal is symmetric and unbounded where the real per-90 rates are skewed and
bounded below, and skew alone lifts the silhouette a partition can reach against this reference,
which is why the copula generator stands beside it.

The copula generator transforms each column of $X$ to normal scores by rank, estimates the
correlation $\hat R$ of those scores, draws $n$ rows from $\mathcal{N}(0, \hat R)$, and maps each
simulated column back through the empirical quantile function of the corresponding observed
column. Every marginal distribution of $X$ is therefore reproduced exactly, including skew,
bounds, ties and any bimodality a single feature carries on its own, and the pairwise rank
correlation is preserved up to sampling noise. What it destroys is joint structure beyond
pairwise rank correlation, which is where two groups differing in several features at once would
live. It tests whether the joint structure carries anything beyond the marginals, not whether
the marginals are unimodal.

The uniform generator rotates $X$ to its principal axes, draws $n$ rows uniformly within the
resulting bounding box, and rotates back. This is the reference distribution of
\citet{tibshirani2001estimating}. It assumes no distributional form at all. Every principal axis
is filled evenly over its full observed range, so a bounded or heavy-tailed axis ranges over
several times its standard deviation and the minor axes carry more spread than in the data,
making the reference more isotropic than the sample it is matched to. That lowers its silhouette
at two clusters, so for the silhouette it is a soft reference; its leading component is flat
rather than peaked, which gives it the largest dip of the three, so for the dip statistic it is
the hard one. Which generator is hardest in each scope is recorded in the metrics file as a
result rather than assumed.

\subsection{Stability, replication and agreement with position}

We draw two hundred bootstrap resamples, refit the clustering on each, and compute the adjusted
Rand index between each bootstrap partition and the reference partition on the players common to
both \citep{hubert1985comparing}, accumulating a consensus matrix that records, for every pair of
players, the fraction of resamples in which the two were assigned together. Clusters whose
members co-assign below 0.6 on average are reported as unstable and described as such rather
than being renamed or merged to make the account tidier.

Replication is tested separately and more stringently, because a resample reuses the same
players. Each season and each league is standardised inside itself and clustered from scratch,
borrowing no centre, no scale and no starting point from the pooled fit. Cluster labels are
arbitrary per fit, so labels are aligned to the pooled partition on the contingency table before
any label-dependent quantity is computed, and an assertion fails the run if an aligned agreement
drops to chance. Three comparisons are reported because they answer different questions: an
independent fit scored against the pooled partition on the same rows; one season's labels scored
against another's on the players present in both, which mixes replication with real change in
those players; and one scope's centroids imported into another scope's players and scored against
that scope's own fit, which isolates replication. Each scope's centroid-difference vector is also
correlated with the pooled one, to check that the scopes recover the same axis rather than merely
a similarly sized partition.

Agreement with the listed positions is quantified by normalised mutual information and by the
adjusted Rand index, against both the position group and the full position string. Players whose
cluster assignment disagrees with their listed group, and players whose Gaussian mixture
posterior is split across components, are extracted as candidate hybrids.

\subsection{Feature-family ablation}

A two-mode result on \ArcFeatures{} features could in principle be carried by one block of them.
The features are grouped into six families, shooting, chance creation, ball progression, pitch
territory, defending and passing accuracy, with the grouping stated in the analysis code so that
it can be audited and asserted to cover every feature exactly once. Each family is removed in
turn and the pre-declared selection rule and the Gaussian null calibration are rerun on what
remains, in every scope. The converse is run as well, each family alone, which shows which
families carry the two-mode structure by themselves. For every ablation the paper records the
number of clusters the rule selects, the standard score of the observed silhouette against the
null, and the adjusted Rand index between the ablated partition and the full one on the same
players, so that a surviving $k$ of two can be distinguished from the same two modes surviving.

\subsection{Recovering and predicting results}

The archive's team tables carry goals and expected goals but not results, so wins, draws and
losses are recovered by summing the goalkeeper table, which credits every match to the
goalkeeper who played it. The recovery is audited: each league-season is checked for the right
number of clubs, every club's credited matches are checked against the schedule with the
exceptions listed, goals for are taken from goals scored while each goalkeeper was on the pitch
and accepted only where the same rows reproduce the goalkeeper table's goals against to within
one, and two published season records are asserted. Points per match, goal difference per match
and league position follow.

Each team-season is then summarised by its minutes-weighted composition on the two-mode partition
and on the six archetypes, using the eligible players, and by the minutes-weighted mean position
of its players along the axis between the two centroids. Prediction is out of sample under three
holdouts, one season at a time, one league at a time and one club at a time, the last being the
strict one since no held-out team-season then shares a club with the training data, by ridge
regression with the penalty chosen by cross-validation inside the training folds, comparing
possession alone, the composition alone and both. The goalkeeper partitions are tested the same
way, against the prediction that a partition sorting keepers by their teams will track points
and a partition describing goalkeeping will not.

\subsection{Archetype profiling and outliers}

For every retained cluster we compute the centroid profile as a standardised deviation from the
mean of its position group, take the eight features with the largest absolute deviation as the
distinguishing features, and list the ten players closest to the centroid and three players near
the boundary. Archetype names are assigned only after inspecting these profiles and the players
they contain, and each name is accompanied by a justification that cites the distinguishing
features and the example players. The names are hypotheses about what the statistical profile
represents, not observations, and they carry no authority beyond the profile that produced them.
Outliers are identified within position group with an isolation forest \citep{liu2008isolation},
reported alongside a Mahalanobis distance so that a distance-based and a partition-based notion
of unusualness can be compared.

\subsection{Supervised validation}

Logistic regression and gradient boosting, the latter with LightGBM \citep{ke2017lightgbm}, are
trained to predict the position group from the feature set under five-fold stratified cross
validation, and we report accuracy, macro F1 and the confusion matrix, every number computed
from out-of-fold predictions. A majority-class baseline runs through the identical folds, so the
reported lift is over the trivial model rather than over zero. SHAP values
\citep{lundberg2017unified} on the boosted model identify which statistics distinguish each
position; the model is fitted on every row and the attributions are computed on a stratified
sub-sample, which is stated with the result. We also record, for each feature and class, the rank
correlation between the feature's value and the attribution it makes, which separates a class
identified by what its members do from one identified by what they do not. Players the classifier
places in the wrong group are cross-referenced against the candidate hybrids from the clustering,
since agreement between two methods that fail in the same place is more informative than either
alone.

A secondary ridge regression predicts non-penalty expected goals per ninety from the non-shooting
features only, so that the fitted value represents the shot volume a player's role would
ordinarily generate and the residual represents the part not explained by role. We report the
cross validated coefficient of determination and treat the residual ranking with heavy caution,
because a residual from a model with modest explanatory power is mostly noise.

\subsection{Dimensionality reduction and the remaining analyses}

Principal component analysis is fitted to the standardised feature set, both globally and
separately within each position group. Uniform manifold approximation and projection
\citep{mcinnes2018umap} is fitted over a grid of twelve configurations, with the number of
neighbours in \{10, 15, 30, 50\} crossed with a minimum distance in \{0.0, 0.1, 0.5\}, and
configurations are ranked by trustworthiness, which measures the degree to which points close in
the embedding were also close in the original space. Selection by an explicit criterion matters
here because a manifold method will produce a visually appealing picture at almost any setting.
t-distributed stochastic neighbour embedding \citep{vandermaaten2008visualizing} is computed at a
perplexity of 30 for visual contrast only, and is never used as an input to clustering, because
its objective preserves local neighbourhoods at the expense of the global geometry a
distance-based clustering algorithm depends upon.

Each club is summarised in two ways, as a minutes-weighted centroid of its outfield players in
the shared principal component space and as the share of its outfield minutes played by each
archetype, and clubs are then clustered hierarchically on archetype composition and compared
against final league position. Player similarity is computed as cosine similarity within position
group. Per-90 rates computed over partial seasons are noisy, and that noise is largest where the
minutes are fewest, so we apply empirical Bayes shrinkage toward the position group mean with the
degree of shrinkage determined by minutes played \citep{efron1977steins}, refit the clustering on
the shrunken features, and report how many assignments change; a partition that dissolves under
shrinkage was substantially a description of sampling noise.
